\chapter{矩阵与图}\label{ch:g12:matrix}

\omterm{def:g12:matrix:matrix}{矩阵}是由数排成的矩形阵列，其加法与乘法的规则经过设计，
使\omterm{def:g12:matrix:matrix}{矩阵}代数能够表示线性
变换的复合。\omterm{def:g12:matrix:matrix}{矩阵}可以求解\omterm{def:g10:lines:system}{线性方程组}，驱动耦合的递推
\omterm{def:g12:seq:sequence}{数列}，并计算网络中的游走数——这正是搜索引擎与
最短路径算法背后的数学。

\section{矩阵代数}

\begin{definition}[矩阵]\label{def:g12:matrix:matrix}
称 \emph{$m \times n$ 矩阵}\index{矩阵}为由\omterm{def:g10:numbers:sets}{实数}构成的表格，它有
$m$ 行 $n$ 列：
$A = (a_{ij})$，其中 $a_{ij}$ 为第 $i$ 行第 $j$ 列的元素。
同型矩阵按对应元素相加，并
规定 $\lambda A = (\lambda a_{ij})$。
\end{definition}

\begin{definition}[矩阵乘积]\label{def:g12:matrix:product}
设 $A$ 为 $m \times n$ \omterm{def:g12:matrix:matrix}{矩阵}，$B$ 为 $n \times p$ \omterm{def:g12:matrix:matrix}{矩阵}，则乘积 $AB$ 是
一个 $m \times p$ \omterm{def:g12:matrix:matrix}{矩阵}，其 $(i,j)$ 元素为
\[
(AB)_{ij} = \sum_{k=1}^{n} a_{ik} b_{kj}
\]
（即“第 $i$ 行（取自 $A$）乘第 $j$ 列（取自 $B$）”的规则）。
\end{definition}

\begin{example}
$\begin{pmatrix} 1 & 2\\ 3 & 4\end{pmatrix}
\begin{pmatrix} 0 & 1\\ 1 & 1\end{pmatrix}
= \begin{pmatrix} 2 & 3\\ 4 & 7\end{pmatrix}$，
\quad 而 \quad
$\begin{pmatrix} 0 & 1\\ 1 & 1\end{pmatrix}
\begin{pmatrix} 1 & 2\\ 3 & 4\end{pmatrix}
= \begin{pmatrix} 3 & 4\\ 4 & 6\end{pmatrix}$：
\emph{\omterm{def:g12:matrix:matrix}{矩阵}乘法不满足交换律}。
\end{example}

\begin{proposition}[矩阵代数的运算规则]\label{prop:g12:matrix:rules}
只要各\omterm{def:g12:matrix:matrix}{矩阵}的大小使乘积有意义，则下式成立：
\[
(AB)C = A(BC), \qquad A(B + C) = AB + AC, \qquad (A+B)C = AC + BC,
\]
以及\emph{单位\omterm{def:g12:matrix:matrix}{矩阵}} $I_n$（对角线元素为 1，其余元素为 0）
满足 $I_m A = A I_n = A$，其中 $A$ 的大小为 $m \times n$。
\end{proposition}

\begin{proof}
以上各式均由
\cref{def:g12:matrix:product} 出发逐元素验证；结合律是唯一需要细究的，
它归结为交换两个有限和的次序：
$\bigl((AB)C\bigr)_{ij} = \sum_l \left(\sum_k a_{ik}b_{kl}\right) c_{lj}
= \sum_k a_{ik} \left(\sum_l b_{kl} c_{lj}\right)
= \bigl(A(BC)\bigr)_{ij}$。
\end{proof}

\begin{definition}[逆矩阵]\label{def:g12:matrix:inverse}
称方阵 $A$（$n$ 阶）为\emph{可逆}\index{矩阵!逆矩阵}，若
存在\omterm{def:g12:matrix:matrix}{矩阵} $B$ 使 $AB = BA = I_n$；此时 $B$ 唯一，记为
$A^{-1}$。
\end{definition}

\begin{proposition}[$2\times2$ 矩阵的逆]\label{prop:g12:matrix:inverse2x2}
设 $A = \begin{pmatrix} a & b\\ c & d\end{pmatrix}$，且
$\det A = ad - bc$（\emph{行列式}\index{行列式}）。则 $A$
\omterm{def:g12:matrix:inverse}{可逆}当且仅当 $\det A \neq 0$，此时
\[
A^{-1} = \frac{1}{ad - bc}\begin{pmatrix} d & -b\\ -c & a\end{pmatrix}.
\]
\end{proposition}

\begin{proof}
直接计算可得
$A \begin{pmatrix} d & -b\\ -c & a\end{pmatrix}
= \begin{pmatrix} d & -b\\ -c & a\end{pmatrix} A = (ad - bc) I_2$；若
$ad - bc \neq 0$，两边除以它即可。反之，若 $ad - bc = 0$，则 $A$ 的各列
成比例，$AB$ 的各列（对任意 $B$）亦然；但
$I_2$ 的各列不成比例，故无 $B$ 能使 $AB = I_2$。
\end{proof}

\begin{method}[线性方程组]\label{met:g12:matrix:systems}
该\omterm{def:g10:algebra:equation}{方程}组即\omterm{def:g12:matrix:matrix}{矩阵}\omterm{def:g10:algebra:equation}{方程}
$\begin{cases} ax + by = e\\ cx + dy = f \end{cases}$
$AX = Y$，其中
$X = \begin{pmatrix} x \\ y\end{pmatrix}$，
$Y = \begin{pmatrix} e \\ f\end{pmatrix}$。若 $\det A \neq 0$，则其唯一
解为 $X = A^{-1}Y$。同一形式亦适用于
$n$ 个\omterm{def:g10:algebra:equation}{方程}、$n$ 个未知数的情形。
\end{method}

\section{矩阵的幂与递推数列}

\begin{definition}\label{def:g12:matrix:power}
对方阵 $A$ 及 $k \in \N$，记 $A^k = A \times \dots \times A$
（共 $k$ 个因子），并规定 $A^0 = I$。
\end{definition}

\begin{method}[数量矩阵加幂零矩阵及可对角化的情形]\label{met:g12:matrix:powers}
计算 $A^k$ 的两种标准方法：
\begin{itemize}
  \item 若 $A = \lambda I + N$，其中 $N^2 = 0$，
        则由二项式定理（因 $I$ 与 $N$ 可交换）化为两项：
        $A^k = \lambda^k I + k \lambda^{k-1} N$。
  \item 若能找到\omterm{def:g12:matrix:inverse}{可逆}\omterm{def:g12:matrix:matrix}{矩阵} $P$ 使 $A = PDP^{-1}$，且 $D$ 为对角\omterm{def:g12:matrix:matrix}{矩阵}，
        则 $A^k = P D^k P^{-1}$，且 $D^k$ 可逐元素计算。
        （求此类 $P$ 的系统方法属于
        \emph{对角化}理论，大学阶段才\omterm{def:g10:algebra:expand}{展开}；在本阶段 $P$ 是
        给定的。）
\end{itemize}
\end{method}

\begin{example}[耦合数列]\label{ex:g12:matrix:coupled}
设 $u_{n+1} = 3u_n + v_n$，$v_{n+1} = u_n + 3v_n$。令
$X_n = \begin{pmatrix} u_n\\ v_n \end{pmatrix}$ 及
$A = \begin{pmatrix} 3 & 1\\ 1 & 3\end{pmatrix}$，则
$X_{n+1} = AX_n$，故 $X_n = A^n X_0$。辅助\omterm{def:g12:seq:sequence}{数列}
$s_n = u_n + v_n$ 与 $d_n = u_n - v_n$ 满足 $s_{n+1} = 4s_n$ 及
$d_{n+1} = 2d_n$，故 $s_n = 4^n s_0$，$d_n = 2^n d_0$，且
\[
u_n = \frac{4^n(u_0+v_0) + 2^n(u_0-v_0)}{2}, \qquad
v_n = \frac{4^n(u_0+v_0) - 2^n(u_0-v_0)}{2}.
\]
（背后的思路：$(1,1)$ 与 $(1,-1)$ 是 $A$ 的特征\omterm{def:g11:vect:vector}{向量}方向。）
\end{example}

\section{图与游走}

\begin{definition}[图、邻接矩阵]\label{def:g12:matrix:graph}
一个\emph{图}\index{图}由顶点 $1, 2, \dots, n$ 及
连接某些顶点对的边组成（对于\emph{有向}
图，边为有序对）。其\emph{邻接矩阵}\index{邻接矩阵}是
$n \times n$ \omterm{def:g12:matrix:matrix}{矩阵} $M$，其中 $m_{ij} = 1$ 当有边从 $i$ 到
$j$，否则为 $0$。长度为 $k$ 的从 $i$ 到 $j$ 的\emph{游走}
是由 $k$ 条相继边组成的、从 $i$ 到 $j$ 的序列。
\end{definition}

\begin{omfigure}
\begin{tikzpicture}[baseline=(current bounding box.center)]
  \node[circle, draw, inner sep=2.5pt] (v1) at (0,0) {$1$};
  \node[circle, draw, inner sep=2.5pt] (v2) at (2.6,1.2) {$2$};
  \node[circle, draw, inner sep=2.5pt] (v3) at (2.6,-1.2) {$3$};
  \draw[-Stealth, omDef, thick] (v1) -- (v2);
  \draw[-Stealth, omDef, thick] (v2) -- (v3);
  \draw[-Stealth, omDef, thick] (v1) to[bend right=18] (v3);
  \draw[-Stealth, omDef, thick] (v3) to[bend right=18] (v1);
\end{tikzpicture}
\qquad\qquad
$M = \begin{pmatrix} 0 & 1 & 1\\ 0 & 0 & 1\\ 1 & 0 & 0 \end{pmatrix}$

{\small 有向图及其\omterm{def:g12:matrix:graph}{邻接矩阵}
（见\cref{exo:g12:matrix:6}）：$m_{ij} = 1$ 当且仅当存在从 $i$ 到
$j$ 的边。}
\end{omfigure}

\begin{theorem}[游走计数]\label{thm:g12:matrix:walks}
长度为 $k$、从顶点 $i$ 到顶点 $j$ 的游走数，
为 $(i,j)$ 位置上的元素，该\omterm{def:g12:matrix:matrix}{矩阵}为 $M^k$。
\end{theorem}

\begin{proof}
对 $k$ 作归纳。当 $k = 1$ 时，即为 $M$ 的定义。假设
对 $k$ 的结论成立。长度为 $k+1$ 的从 $i$ 到 $j$ 的游走，是由长度为 $k$ 的、从 $i$ 到某顶点 $l$ 的游走
后接一条从 $l$ 到 $j$ 的边组成，由
加法原理与乘法原理，其数目为
\[
\sum_{l=1}^{n} \bigl(M^k\bigr)_{il}\, m_{lj} = \bigl(M^{k+1}\bigr)_{ij}.
\qedhere
\]
\end{proof}

\begin{example}
对三角形图（$3$ 个顶点，每两个顶点都相连），
$M = \begin{pmatrix} 0&1&1\\ 1&0&1\\ 1&1&0\end{pmatrix}$ 及
$M^2 = \begin{pmatrix} 2&1&1\\ 1&2&1\\ 1&1&2\end{pmatrix}$：从每个顶点出发，
回到自身的长度为 $2$ 的游走有 $2$ 条（经由任一邻点），
到其余每个顶点的游走各有 $1$ 条。
\end{example}

\section{练习}

\begin{exercise}[$\star$]\label{exo:g12:matrix:1}
设 $A = \begin{pmatrix} 1 & 2\\ 0 & 1 \end{pmatrix}$，
$B = \begin{pmatrix} 2 & 0\\ 1 & 1 \end{pmatrix}$。计算 $A + B$，$AB$，
$BA$ 及 $A^2$。
\end{exercise}

\begin{exercise}[$\star$]\label{exo:g12:matrix:2}
判断下列\omterm{def:g12:matrix:matrix}{矩阵}是否\omterm{def:g12:matrix:inverse}{可逆}，并求出
\omterm{def:g12:matrix:inverse}{可逆}时的逆\omterm{def:g12:matrix:matrix}{矩阵}：
\[
A = \begin{pmatrix} 2 & 5\\ 1 & 3\end{pmatrix}, \qquad
B = \begin{pmatrix} 3 & 6\\ 2 & 4\end{pmatrix}.
\]
\end{exercise}

\begin{exercise}[$\star$]\label{exo:g12:matrix:3}
用\omterm{def:g12:matrix:matrix}{矩阵}求逆法解\omterm{def:g10:algebra:equation}{方程}组：
$\begin{cases} 2x + 5y = 1\\ x + 3y = 2 . \end{cases}$
\end{exercise}

\begin{exercise}[$\star\star$]\label{exo:g12:matrix:4}
设 $A = \begin{pmatrix} 2 & 1\\ 0 & 2\end{pmatrix} = 2I + N$，其中
$N = \begin{pmatrix} 0 & 1\\ 0 & 0 \end{pmatrix}$.
\begin{enumerate}
  \item 验证 $N^2 = 0$，且 $I$ 与 $N$ 可交换。
  \item 由此求出 $A^k$（对所有 $k \in \N$），并
        用直接计算验证 $k=2$ 时的公式。
\end{enumerate}
\end{exercise}

\begin{exercise}[$\star\star$]\label{exo:g12:matrix:5}
设 $A = \begin{pmatrix} 0 & 1\\ 1 & 1\end{pmatrix}$，记 $F_n$ 为斐波那契
\omterm{def:g12:seq:sequence}{数列}（$F_0 = 0$，$F_1 = 1$，$F_{n+2} = F_{n+1} + F_n$）。用归纳法证明：
对 $n \geq 1$，
\[
A^n = \begin{pmatrix} F_{n-1} & F_n\\ F_n & F_{n+1}\end{pmatrix},
\]
并由此推出恒等式 $F_{n+1}F_{n-1} - F_n^2 = (-1)^n$。
（提示：\omterm{def:g11:vect:det}{行列式}满足乘法法则 $\det(MN) = \det M \det N$，你可以
对 $2\times2$ \omterm{def:g12:matrix:matrix}{矩阵}加以验证。）
\end{exercise}

\begin{exercise}[$\star\star$]\label{exo:g12:matrix:6}
顶点为 $\{1, 2, 3\}$ 的有向图，其边为
$1\to2$、$2\to3$、$3\to1$ 和 $1\to3$。
\begin{enumerate}
  \item 写出\omterm{def:g12:matrix:graph}{邻接矩阵} $M$，并计算 $M^2$ 和 $M^3$。
  \item 长度为 $3$ 的游走中，从 $1$ 到 $1$ 的有多少条？请将它们全部列出。
\end{enumerate}
\end{exercise}

\begin{exercise}[$\star\star$]\label{exo:g12:matrix:7}
某汽车共享公司在两座城市 $A$ 与 $B$ 之间调配车辆。每
周，$80\%$ 的 $A$ 市车辆留在 $A$ 市，$20\%$ 迁往 $B$ 市；
$30\%$ 的 $B$ 市车辆迁往 $A$ 市，$70\%$ 留在原城市。设 $a_n, b_n$ 分别为
各城市车队所占的比例。
\begin{omfigure}
\begin{tikzpicture}
  \node[circle, draw, inner sep=3pt] (A) at (0,0) {$A$};
  \node[circle, draw, inner sep=3pt] (B) at (3.2,0) {$B$};
  \draw[-Stealth, omDef, thick] (A) to[bend left=25]
    node[above, font=\small] {$0.2$} (B);
  \draw[-Stealth, omDef, thick] (B) to[bend left=25]
    node[below, font=\small] {$0.3$} (A);
  \draw[-Stealth, omDef, thick] (A) to[loop left]
    node[left, font=\small] {$0.8$} (A);
  \draw[-Stealth, omDef, thick] (B) to[loop right]
    node[right, font=\small] {$0.7$} (B);
\end{tikzpicture}
\end{omfigure}
\begin{enumerate}
  \item 写出 $X_{n+1} = MX_n$，其中
        $X_n = \begin{pmatrix} a_n\\ b_n\end{pmatrix}$，并确定 $M$。
  \item 求平衡比例（解 $MX = X$，
        并满足 $a + b = 1$）。
  \item 证明 $c_n = a_n - 0.6$ 满足 $c_{n+1} = 0.5\,c_n$，并
        由此推出车队\omterm{def:g11:prob:rv}{分布}\omterm{def:g12:seq:limit}{收敛}于平衡状态。
\end{enumerate}
\end{exercise}

\begin{exercise}[$\star\star\star$]\label{exo:g12:matrix:8}
设 $A = \begin{pmatrix} 3 & 1\\ 1 & 3 \end{pmatrix}$，
$P = \begin{pmatrix} 1 & 1\\ 1 & -1 \end{pmatrix}$。
\begin{enumerate}
  \item 计算 $P^{-1}$，再求 $D = P^{-1}AP$，并验证 $D$ 是
        对角\omterm{def:g12:matrix:matrix}{矩阵}。
  \item 由此求 $A^n$ 的显式表达式，并与
        \cref{ex:g12:matrix:coupled} 中的结果比较。
\end{enumerate}
\end{exercise}

\section{问题：认识斐波那契数列的矩阵（以及
天气）}

\begin{problem}[{周末问题——一个 $2 \times 2$ 矩阵
包含全部斐波那契数，一个马尔可夫矩阵预报长期
天气，一个特征向量价值十亿美元}]
\label{pb:g12:matrix:1}
\omterm{def:g12:matrix:matrix}{矩阵}是这样一台机器：输入一个状态，返回下一个状态
——因而它的\emph{幂}蕴含整个未来。本题
从一个惊人的\omterm{def:g12:matrix:matrix}{矩阵}开始，其幂列出
斐波那契数（且每个恒等式都只需一行即可证明），
接着把天气视为马尔可夫链，求其稳态，并
最后以搜索引擎赖以构建的特征\omterm{def:g11:vect:vector}{向量}作结
（见\cref{thm:g12:matrix:walks}，
\cref{met:g12:matrix:powers}）。

\medskip
\noindent\textbf{第一部分——熟练运算。}
\begin{enumerate}
  \item 令 $A = \begin{pmatrix} 1 & 2\\ 3 & 4\end{pmatrix}$
        及 $B = \begin{pmatrix} 0 & 1\\ 1 & 0\end{pmatrix}$：
        计算 $AB$ 与 $BA$，并判断乘法是否满足交换律？
  \item 求 $\begin{pmatrix} 2 & 1\\ 5 & 3\end{pmatrix}$ 的逆\omterm{def:g12:matrix:matrix}{矩阵}
        （见\cref{prop:g12:matrix:inverse2x2}），并用逆\omterm{def:g12:matrix:matrix}{矩阵}
        解\omterm{def:g10:algebra:equation}{方程}组 $2x + y = 4$，$5x + 3y = 7$。
  \item 设 $N = \begin{pmatrix} 0 & 1\\ 0 & 0\end{pmatrix}$，
        计算 $N^2$，并由此推出
        $(I + N)^n = I + nN$ 对所有 $n$ 成立。
  \item 三角形图（三个顶点两两相连）：
        写出其\omterm{def:g12:matrix:graph}{邻接矩阵} $A$，计算 $A^3$，并
        解释其对角线上元素的含义
        （见\cref{thm:g12:matrix:walks}）。
  \item 对 $D = \begin{pmatrix} 2 & 0\\ 0 & \frac12
        \end{pmatrix}$：求 $D^n$，并说明
        $n \to \infty$ 时的行为。
\end{enumerate}

\medskip
\noindent\textbf{第二部分——斐波那契\omterm{def:g12:matrix:matrix}{矩阵}。}
设 $F = \begin{pmatrix} 1 & 1\\ 1 & 0\end{pmatrix}$，并设
$F_1 = F_2 = 1, F_3 = 2, \dots$ 是
\cref{pb:g11:seq:1}中的斐波那契数。
\begin{enumerate}[resume]
  \item 计算 $F^2$、$F^3$、$F^4$，并猜想 $F^n$ 的一般形式，
        该形式用斐波那契数表示。
  \item 用归纳法证明猜想
        $F^n = \begin{pmatrix} F_{n+1} & F_n\\ F_n & F_{n-1}
        \end{pmatrix}$。
  \item 对等式两边取\omterm{def:g11:vect:det}{行列式}（\omterm{def:g12:matrix:matrix}{矩阵}乘积的\omterm{def:g11:vect:det}{行列式}
        等于各自\omterm{def:g11:vect:det}{行列式}之积，若从未见过，可
        在 $2 \times 2$ \omterm{def:g12:matrix:matrix}{矩阵}上验证）：由此推出
        \emph{卡西尼恒等式}
        $F_{n+1}F_{n-1} - F_n^2 = (-1)^n$——这是
        “平方项相消”的核心，一行即可证明。
  \item 由 $F^{m+n} = F^m F^n$，读出右上角元素，
        并推出\emph{加法公式}
        \[
        F_{m+n} = F_{m+1} F_n + F_m F_{n-1} .
        \]
        在 $m = n = 3$ 时验证该式。
  \item 由加法公式（对 $k$ 用归纳法）推出
        $F_n$ \omterm{def:g12:arith:divides}{整除} $F_{kn}$，并验证
        $F_3 \mid F_6$ 与 $F_3 \mid F_9$。
  \item 要计算 $F_{100}$，不必做 $100$ 次\omterm{def:g12:matrix:matrix}{矩阵}乘法，
        只需反复平方（$F^2, F^4, F^8, \dots$），
        再将其组合。多少次\omterm{def:g12:matrix:matrix}{矩阵}乘法就足够，
        这是初中分册中哪一种古老的乘法技巧，
        被推广到\omterm{def:g12:matrix:matrix}{矩阵}上？
\end{enumerate}

\medskip
\noindent\textbf{第三部分——天气机器。}
某城市中：晴天之后，次日为晴天的
\omterm{def:g10:proba:distribution}{概率}为 $0.8$；雨天之后，次日为晴天的\omterm{def:g10:proba:distribution}{概率}为
$0.4$。将每天的\omterm{def:g11:prob:rv}{分布}记为一列
$\binom{p_{\text{晴}}}{p_{\text{雨}}}$，演化由下式给出。
\[
M = \begin{pmatrix} 0.8 & 0.4\\ 0.2 & 0.6 \end{pmatrix}.
\]
\begin{enumerate}[resume]
  \item 验证 $M$ 的每一列之和为 $1$，并说明
        为何任何天气机器都必须具有此性质。
  \item 今天是晴天。求明天与后天的
        天气预报。
  \item 求\emph{稳态}：\omterm{def:g11:prob:rv}{分布} $v$，
        满足 $Mv = v$（分量之和为 $1$）。
        长期来看，晴天占多大比例？
  \item 从雨天出发，取初始\omterm{def:g11:prob:rv}{分布} $\binom01$，\omterm{def:g12:limcont:continuity}{连续}作用 $M$ 四次，
        并在每一步记录其与稳态的距离。
        每一步间距缩小的倍数是多少？
        这属于哪一种\omterm{def:g12:seq:limit}{收敛}？
  \item 迷你版 PageRank：三个网页，链接关系为
        $A \to B$，$A \to C$，$B \to C$，$C \to A$。随机浏览者
        以均匀随机的方式沿一条出链前进。
        写出转移\omterm{def:g12:matrix:matrix}{矩阵}，求稳态，
        并对网页进行排序。
  \item 解释该排序：为何 $C$ 的得分与
        $A$ 一样高，尽管其接收的链接来自更少的网页？
        稳态究竟度量的是什么？（真实的 PageRank
        对死端与跳转还需加入一个阻尼因子；
        特征\omterm{def:g11:vect:vector}{向量}的思想正是这一点。）
\end{enumerate}

\medskip
\noindent\textbf{第四部分——对角化的红利。}
\begin{enumerate}[resume]
  \item 两个相互耦合的量满足
        $u_{n+1} = 3u_n + v_n$，$v_{n+1} = u_n + 3v_n$，即
        \cref{exo:g12:matrix:8}中的\omterm{def:g12:matrix:matrix}{矩阵} $A$。利用该题
        的对角化，
        （$D = \operatorname{diag}(4, 2)$），给出 $u_n$ 的闭式表达，
        其中 $u_0 = 1$、$v_0 = 0$，并
        与 $n = 1, 2, 3$ 时的直接计算结果相核对。
  \item 用一两句话说明：对角化对耦合系统
        究竟\emph{做了什么}？以及
        第 14 题中的马尔可夫稳态，在何种意义上
        也是一个特征\omterm{def:g11:vect:vector}{向量}的问题？
  \item 终章——本周末\omterm{def:g12:matrix:matrix}{矩阵}的三副面孔：
        记账（\omterm{def:g10:algebra:equation}{方程}组与逆）、组合计数
        （以幂计数的路径与链接），以及演化
        （斐波那契\omterm{def:g12:seq:sequence}{数列}、天气与网页——未来由特征方向读出），
        每项各用一句话概括，并给出
        前瞻：大学分册中的线性代数
        将使上述每一面都成为一门理论。
\end{enumerate}
\end{problem}
